\documentclass[10pt,a4paper]{amsart}
\usepackage[margin=19mm]{geometry}
\usepackage{amsmath,amssymb,mathtools}
\usepackage[scr=rsfs,cal=euler]{mathalfa}
\usepackage{xfrac}
\usepackage[unicode,hidelinks,bookmarks=false]{hyperref}
\newcommand{\ie}{i.e.\ }
\newcommand{\cf}{cf.\ }
\newcommand{\resp}{respectively\ }
\newcommand{\Subsection}{\subsection}
\providecommand{\nobreakdash}{\nobreak\hbox{-}}
\setlength{\emergencystretch}{2em}
\multlinegap0pt
\usepackage{lmodern}
\usepackage{mathtools}
\usepackage{enumitem}
\allowdisplaybreaks
\newtheorem{theorem}{Theorem}[section]
\newtheorem{proposition}[theorem]{Proposition}
\newtheorem{corollary}[theorem]{Corollary}
\newtheorem{lemma}[theorem]{Lemma}
\newtheorem{definition}[theorem]{Definition}
\newtheorem{example}[theorem]{Example}
\newtheorem{remark}[theorem]{Remark}
\newcommand{\kk}{\Bbbk}
\newcommand{\Der}{\operatorname{Der}}
\newcommand{\Sym}{\operatorname{Sym}}
\newcommand{\id}{\operatorname{id}}
\newcommand{\tr}{\operatorname{tr}}
\newcommand{\Pcal}{\mathcal P}
\newcommand{\Fcot}{F_{\mathrm{cot}}}
\begin{document}
\title[From Lie--Rinehart Algebras to F-Manifold Algebras]{From Lie--Rinehart Algebras to $F$-Manifold Algebras}
\author{Yufeng Pei; Yunhe Sheng}
\date{}
\maketitle
\noindent\emph{Source and licence.} From Lie--Rinehart Algebras to $F$-Manifold Algebras. Version arXiv:2608.12802v1 (CC BY 4.0). This material is distributed under the Creative Commons Attribution 4.0 International licence, http://creativecommons.org/licenses/by/4.0/, as recorded in the arXiv OAI licence metadata for this exact version and shown on the abstract page https://arxiv.org/abs/2608.12802v1. Full rights notice: licenses/arxiv-2608.12802-CC-BY-4.0.txt.\par\smallskip
\noindent\textit{Editorial scope.} Counted unit: the original \texttt{theorem} environment \texttt{thm:first-order-truncation} at source lines 436--447 together with its complete proof at lines 449--476; the delivered document keeps source lines 247--477 so that every referenced label and every citation resolves inside it. Native editable source is the paper's own concatenated TeX; the AMS wrapper is editorial formatting only. No publisher class, upstream script or unrelated artwork is redistributed.
	\begin{equation}\label{eq:F-product}
		(a+X)\bullet(b+Y)=ab+aY+bX,
		\qquad a,b\in A,\quad X,Y\in L,
	\end{equation}
	Thus $A$ is a subalgebra of $P$, $L$ is an ideal with
	$L\bullet L=0$, and $a\bullet X=aX$ for $a\in A$ and $X\in L$.
	Define a bracket on $P$ by
	\begin{equation}\label{eq:F-bracket}
		[a+X,b+Y]
		=\rho(X)(b)-\rho(Y)(a)+[X,Y]_L,
		\qquad a,b\in A,\quad X,Y\in L.
	\end{equation}
	Equivalently,
	\[
	[A,A]=0,
	\qquad
	[X,a]=\rho(X)(a),
	\qquad
	[X,Y]=[X,Y]_L.
	\]
	This is the semidirect bracket determined by the anchor action on \(A\) and
	the Lie bracket on \(L\).
	
	
	
	We compute the Leibnizator of the bracket with respect to $\bullet$.
	
	\begin{lemma}\label{lem:Leibnizator-formula}
		For $a+X, b+Y, c+Z\in A\oplus L$, we have
		\begin{equation}\label{eq:Leibnizator-formula}
			\Pcal_{a+X}(b+Y,c+Z)=\rho(Y)(a)Z+\rho(Z)(a)Y.
		\end{equation}
		
	\end{lemma}
	
	\begin{proof}
		We expand the three terms in the definition of $\Pcal$.  First,
		\[
		(b+Y)\bullet(c+Z)=bc+bZ+cY.
		\]
		Using \eqref{eq:F-bracket}, the derivation property of $\rho(X)$, the
		$A$-linearity of the anchor, and the Lie--Rinehart identity, we get
		\begin{align*}
			[a+X,(b+Y)\bullet(c+Z)]
			&= [a+X,bc+bZ+cY] \\
			&=\rho(X)(bc)-\rho(bZ+cY)(a)+[X,bZ+cY]_L \\
			&=\rho(X)(b)c+b\rho(X)(c)-b\rho(Z)(a)-c\rho(Y)(a) \\
			&\quad +b[X,Z]_L+c[X,Y]_L+\rho(X)(b)Z+\rho(X)(c)Y.
		\end{align*}
		Second,
		\begin{align*}
			[a+X,b+Y]\bullet(c+Z)
			&=\bigl(\rho(X)(b)-\rho(Y)(a)+[X,Y]_L\bigr)\bullet(c+Z) \\
			&=c\rho(X)(b)-c\rho(Y)(a)+\rho(X)(b)Z-\rho(Y)(a)Z+c[X,Y]_L.
		\end{align*}
		Third,
		\begin{align*}
			(b+Y)\bullet[a+X,c+Z]
			&=(b+Y)\bullet\bigl(\rho(X)(c)-\rho(Z)(a)+[X,Z]_L\bigr) \\
			&=b\rho(X)(c)-b\rho(Z)(a)+\rho(X)(c)Y-\rho(Z)(a)Y+b[X,Z]_L.
		\end{align*}
		Subtracting the last two displayed expressions from the first cancels all
		terms except
		\[
		\rho(Y)(a)Z+\rho(Z)(a)Y.
		\]
		This proves \eqref{eq:Leibnizator-formula}.
	\end{proof}
	
	\begin{theorem}\label{thm:LR-to-FMan}
		Let $(A,L)$ be a Lie--Rinehart algebra.  Then
		$(A\oplus L,\bullet,[-,-])$ is an $F$-manifold algebra.
	\end{theorem}
	
	\begin{proof}
		By \eqref{eq:F-product}, the product $\bullet$ is commutative and associative.
		The bracket is the semidirect product bracket
		for the action $\rho:L\to\Der_\kk(A)$.  It is therefore a Lie bracket: the
		Jacobi identity on $L$ is inherited from $[-,-]_L$, the identity involving two
		elements of $L$ and one element of $A$ follows because $\rho$ is a Lie algebra
		homomorphism, and all remaining components vanish.  It remains to prove the
		Hertling--Manin identity
		\begin{equation}\label{eq:HM-proof-target}
			\Pcal_{(a+X)\bullet (b+Y)}(c+Z,d+W)
			=(a+X)\bullet\Pcal_{b+Y}(c+Z,d+W)+(b+Y)\bullet\Pcal_{a+X}(c+Z,d+W).
		\end{equation}
		
		By Lemma \ref{lem:Leibnizator-formula}, we have
		\begin{align*}
			\Pcal_{(a+X)\bullet (b+Y)}(c+Z,d+W)
			&=\rho(Z)(ab)W+\rho(W)(ab)Z \\
			&=\bigl(a\rho(Z)(b)+b\rho(Z)(a)\bigr)W
			+\bigl(a\rho(W)(b)+b\rho(W)(a)\bigr)Z,
		\end{align*}
		where we used that $\rho(Z)$ and $\rho(W)$ are derivations of $A$.
		
		On the other hand,
		\[
		\Pcal_{b+Y}(c+Z,d+W)=\rho(Z)(b)W+\rho(W)(b)Z\in L,
		\]
		and similarly
		\[
		\Pcal_{a+X}(c+Z,d+W)=\rho(Z)(a)W+\rho(W)(a)Z\in L.
		\]
		Since an element $a+X\in P$ acts by the product $\bullet$ on $L$ only through
		its $A$-component, we have
		\begin{align*}
			(a+X)\bullet\Pcal_{b+Y}(c+Z,d+W)
			&=a\rho(Z)(b)W+a\rho(W)(b)Z,\\
			(b+Y)\bullet\Pcal_{a+X}(c+Z,d+W)
			&=b\rho(Z)(a)W+b\rho(W)(a)Z.
		\end{align*}
		Adding the two equalities gives the expression for
		$\Pcal_{(a+X)\bullet (b+Y)}(c+Z,d+W)$.  Hence \eqref{eq:HM-proof-target} holds.
	\end{proof}
	
	
	
	
	The following corollary characterizes the Poisson case.
	
	\begin{corollary}\label{cor:poisson-specialization}
		The associated $F$-manifold algebra $(A\oplus L,\bullet,[-,-])$ is a
		Poisson algebra if and only if
		\begin{equation}\label{eq:Poisson-condition}
			\rho(Y)(a)Z+\rho(Z)(a)Y=0
		\end{equation}
		for all $a\in A$ and $Y,Z\in L$.
	\end{corollary}
	
	\begin{proof}
		For a commutative associative algebra with a Lie bracket, the Poisson identity
		is equivalent to the vanishing of the Leibnizator.  The assertion therefore
		follows directly from Lemma \ref{lem:Leibnizator-formula}.
	\end{proof}
	
	\begin{remark}\label{rem:LV-poisson-claim}
		It is asserted in \cite[Proposition 13.3.26]{LodayVallette} that $A\oplus L$,
		with the product \eqref{eq:F-product} and bracket \eqref{eq:F-bracket}, is a
		Poisson algebra for an arbitrary Lie--Rinehart algebra.  Corollary
		\ref{cor:poisson-specialization} shows that the additional condition
		\eqref{eq:Poisson-condition} is required.  Theorem
		\ref{thm:LR-to-FMan} gives the valid conclusion without this condition.
	\end{remark}
	
	
	
	
	
	\section{Poisson truncations and the Lie--Rinehart differential}
	\label{sec:ideal}
	
	We now compare this construction with truncations of the canonical
	Poisson algebra associated with a Lie--Rinehart algebra.  This comparison
	also identifies the higher-order obstruction to obtaining a Poisson
	quotient.
	
	Let $(A,L)$ be a Lie--Rinehart algebra and write
	\[
	\Sym_A(L)=\bigoplus_{n\geq0}\Sym_A^n(L).
	\]
	We use the following standard linear Poisson construction; see, for example,
	\cite{Huebschmann1990,Huebschmann1998}.
	
	\begin{proposition}\label{prop:linear-poisson-prelim}
		The symmetric algebra $\Sym_A(L)$ carries a unique Poisson bracket
		determined on generators by
		\begin{equation}\label{eq:linear-poisson}
			\{a,b\}=0,
			\qquad
			\{X,a\}=\rho(X)(a),
			\qquad
			\{X,Y\}=[X,Y]_L,
		\end{equation}
		for $a,b\in A$ and $X,Y\in L$.  Moreover,
		\[
		\{\Sym_A^p(L),\Sym_A^q(L)\}\subseteq \Sym_A^{p+q-1}(L).
		\]
		Here $\Sym_A^{-1}(L)=0$.
	\end{proposition}
	
	\subsection{First-order and higher truncations}
	
	Let
	\[
	I=\bigoplus_{n\geq1}\Sym_A^n(L)
	\]
	be the positive-degree ideal of $\Sym_A(L)$.
	

	\begin{theorem}\label{thm:first-order-truncation}
		As commutative algebras,
		\[
		\Sym_A(L)/I^2\simeq A\oplus L,
		\]
		with the product \eqref{eq:F-product}.  The ideal $I^2$ is a Poisson ideal of
		$\Sym_A(L)$ if and only if
		\[
		\Pcal_a(X,Y)=\rho(Y)(a)X+\rho(X)(a)Y=0
		\]
		for all $a\in A$ and $X,Y\in L$.
	\end{theorem}

	\begin{proof}
		The quotient statement follows from the symmetric grading:
		\[
		\Sym_A(L)/I^2\simeq \Sym_A^0(L)\oplus \Sym_A^1(L)=A\oplus L.
		\]
		The induced multiplication is \eqref{eq:F-product}.
		
		For $X,Y\in L$ and $a\in A$, the Poisson Leibniz rule in $\Sym_A(L)$ gives
		\[
		\{XY,a\}=X\{Y,a\}+Y\{X,a\}
		=\rho(Y)(a)X+\rho(X)(a)Y.
		\]
		By Lemma \ref{lem:Leibnizator-formula}, the same expression is
		$\Pcal_a(X,Y)$ in the $F$-manifold algebra $A\oplus L$.  It has
		degree one.  Hence a nonzero value gives
		$\{I^2,\Sym_A(L)\}\not\subseteq I^2$.
		
		Conversely, assume that the displayed expression vanishes for all $a,X,Y$.
		Since the linear Poisson bracket has  degree $-1$, it is enough, by
		the biderivation property, to check brackets of degree-two monomials with
		algebra generators of degrees $0$ and $1$.  The degree-zero case is the
		assumption.  For $Z\in L$,
		\[
		\{XY,Z\}=X\{Y,Z\}+Y\{X,Z\}=X[Y,Z]_L+Y[X,Z]_L\in I^2.
		\]
		Since $I^2$ is generated by degree-two monomials, this proves
		$\{I^2,\Sym_A(L)\}\subseteq I^2$.
	\end{proof}
	

\begin{thebibliography}{99}
\bibitem[Huebschmann1990]{Huebschmann1990} J.~Huebschmann, \emph{Poisson cohomology and quantization}, \emph{J. Reine Angew. Math.} \textbf{408} (1990), 57--113.
\bibitem[Huebschmann1998]{Huebschmann1998} J.~Huebschmann, \emph{Lie--Rinehart algebras, Gerstenhaber algebras, and Batalin--Vilkovisky algebras}, \emph{Ann. Inst. Fourier (Grenoble)} \textbf{48} (1998), no.~2, 425--440.
\bibitem[LodayVallette]{LodayVallette} J.-L.~Loday and B.~Vallette, \emph{Algebraic Operads}, Grundlehren der mathematischen Wissenschaften, vol.~346, Springer, Heidelberg, 2012.
\end{thebibliography}
\end{document}
